\section*{Discussion}

\subsection*{What media-only sampling can and cannot tell us}

The central result of this paper is exact and independent of any particular drug or platform: a minimal two-compartment liver-chip model, observed through media-only sampling, is structurally unidentifiable by precisely one degree of freedom, and the entire one-parameter family of equally consistent parameter quadruples can be written in closed form (Eq.~\eqref{eq:family}). This is not a statement about insufficient data density or duration --- the three identifiable combinations $a_{11}$, $a_{22}$, $P$ are exactly what the media concentration-time curve can ever determine, however many samples are taken or however long the incubation runs. The practical corollary is that a cell-free non-specific-binding control, already standard practice on some but not all liver-chip platforms, is not merely good laboratory hygiene: it is a mathematical precondition for the remaining kinetic parameters to be identifiable at all from this type of experiment.

\subsection*{A methodological finding: local diagnostics can mislead near a structural degeneracy}

Beyond the structural result, this work surfaced a second, more general methodological finding that we believe is of independent interest. Near a structural (or near-structural) degeneracy, the standard toolkit for quantifying how well a parameter is constrained by data can itself become unreliable. We observed this directly: a Fisher-information covariance approximation gave condition numbers of $10^{12}$ to $10^{17}$, and switching between a plain matrix inverse and a pseudoinverse on identical fitted data changed the estimated relative standard error of $CL_{int}$ by two to three orders of magnitude. Our first response --- adopting the pseudoinverse as the more numerically stable choice --- turned out to be informative in the wrong direction: a direct profile-likelihood check showed that the pseudoinverse-based estimate was misleadingly optimistic relative to the genuinely flat likelihood surface. A single-start profile likelihood, often treated as the more rigorous alternative to a covariance approximation, was itself unreliable here, producing a non-monotonic profile attributable to the optimizer failing to converge at some grid points. Only a multi-start profile likelihood, which is substantially more computationally expensive, gave a profile we are confident reflects the actual information content of the data. We report this progression explicitly, including the dead ends, because it is a cautionary tale with a clear methodological implication: near a known or suspected structural degeneracy, point estimates and standard practical-identifiability diagnostics should not be trusted without a multi-start or otherwise global check.

\subsection*{Structural identifiability is necessary, not sufficient}

The multi-start profile likelihood result in Fig.~\ref{fig:profile}b is, on its own, a finding that qualifies the paper's main structural message. Resolving the exact structural degeneracy by fixing $k_b$ and evaporation $e$ at their true values is, by Eq.~\eqref{eq:family}, sufficient to make the model structurally identifiable --- a unique solution exists in the noise-free, infinite-data limit. It did not, in the regime we examined (low $CL_{int}$ relative to $CL_d$), translate into a practically well-constrained estimate of $CL_{int}$ under either sparse or dense realistic sampling. This is consistent with the broader point made by Milani et al.\ \citep{milani2024} that structural identifiability does not guarantee precise practical estimation, and it sharpens that point by showing a concrete case in which the gap persists even after the specific structural obstruction (the $k_b$-parametrized equivalence family) has been removed. We did not characterize how this result depends on the regime (the ratio of $CL_{int}$ to $CL_d$, captured by the dimensionless group $\mu$ in Eq.~\eqref{eq:nondim}); this is a natural direction for further work.

\subsection*{Relationship to prior identifiability work on organ-on-a-chip systems}

Docci et al.\ \citep{docci2022} performed a priori structural identifiability analysis, via DAISY, for three candidate models of their specific PhysioMimix configuration. Milani et al.\ \citep{milani2022} assessed structural and practical identifiability for a mycophenolate mofetil gut-liver system, and Milani et al.\ \citep{milani2024} generalized this into an a priori experimental-design workflow (structural identifiability, global dynamic sensitivity analysis, and simulation-based verification) demonstrated on midazolam. The distinction between that body of work and this paper is categorical, not incremental: each of those analyses asks whether \emph{one} model is identifiable, for \emph{one} drug, on \emph{one} device, and answers that question by running the check on that specific case. This paper instead proves a theorem about the minimal model class itself --- exact, closed-form, and true for every drug and every platform sharing that structure, before any specific experiment is designed or any specific DAISY check is run. A reader of Docci et al.\ or Milani et al.\ learns whether \emph{that} experiment was identifiable; a reader of this paper learns what any media-only liver-chip experiment of this structure can and cannot claim, in advance. Our contribution is complementary to, not a replacement for, model-specific checks: it tells a reader what to expect before running one, together with a direct demonstration that the practical-identifiability side of any such workflow requires a global (multi-start or equivalent) method to be trusted. The three-compartment DigiLoCS digital twin \citep{aravindakshan2025} adds a further compartment (media, interstitium, and intracellular space) to the structure considered here; whether an analogous equivalence family exists in that richer model, and whether it is removed by the same NSB-control strategy, is an open question raised directly by our result.

\subsection*{Toward a minimum reporting standard}

Existing microphysiological system quality-management recommendations \citep{pamies2024} address biological reproducibility, not what a study must report for independent identifiability assessment. Table~\ref{tab:checklist} sets out items we believe a liver-chip PK study should report if its CLint estimate is intended to support IVIVE: the first eight are items that should simply be reported; the final two --- a structural identifiability assessment for the specific observation scheme used, and a practical-identifiability or uncertainty analysis robust to the kind of degeneracy demonstrated here --- are analyses we recommend performing, not merely data to disclose.

\subsection*{Limitations}

$K_p$, $CL_d$, and the specific numerical values used in our worked examples are illustrative and not fitted to any specific compound; the qualitative structural and methodological results do not depend on these choices, but their magnitudes do. The closed-form structural result (Eq.~\eqref{eq:family}) is exact for the constant-volume reduction of the model; we did not perform a DAISY-based or equivalent structural analysis of the complete five-parameter, time-varying-volume model that includes evaporation as a jointly fitted (rather than externally measured) quantity. We examined the practical-identifiability question in one illustrative regime (low $CL_{int}$ relative to $CL_d$, CN Bio hardware); we did not exhaustively map how the multi-start profile likelihood result varies across the full dimensionless parameter space defined by Eq.~\eqref{eq:nondim}, nor did we test it against other commercial liver-chip platforms.

A preliminary investigation into this last point is worth reporting explicitly, as a direction for future work rather than a settled result. Holding $\mu$ and $\beta$ fixed and varying only $r = K_pV_c/V_m$ (the cell-compartment volume relative to media, a hardware-determined quantity, $r\approx0.009$ for the CN Bio geometry used throughout this paper), a zero-noise diagnostic showed that the sensitivity of the model fit to an incorrect $CL_{int}$ increases sharply with $r$: a 100-fold increase in $V_c$ increased the ratio of fit degradation (at a 10-fold wrong $CL_{int}$) to the fit quality (at the true value) by roughly eleven orders of magnitude. This suggests $r$, not only $\mu$ and $\beta$, governs practical identifiability, and that platforms with a larger cell-compartment volume relative to media may be intrinsically better powered to estimate $CL_{int}$ from media-only data, independent of sampling design. We were not able, within the scope of this paper, to translate this into a controlled comparison across $r$ under realistic noise, because changing $V_c$ also changes the system's intrinsic timescales (the exchange rate $CL_d/V_c$ depends on $V_c$ directly), so a fixed sampling schedule cannot be reused across different $r$ without confounding the comparison; a proper test requires re-deriving an appropriately rescaled sampling design for each $r$. We flag this explicitly as an open, well-defined question for follow-up work, rather than including an inconclusive figure here. We make no claim to reproduce any specific published dataset's bias magnitude, nor that our reduced model represents every liver-chip analysis in current use.
